\textbf{Phần 1. Trắc nghiệm lựa chọn}
\setcounterpageref{ex}{0}

\begin{ex} %[2D8V38]
Trong phép nhân $5 \times 4 = 20$, số $20$ được gọi là
\choice
{Thừa số}
{\True Tích}
{Số bị chia}
{Tổng}
\loigiai{
20 là kết quả của phép nhân nên được gọi là Tích.
\textbf{Đáp án của câu hỏi là B}
}
\end{ex}

\begin{ex} %[2D8V38]
Biết thừa số thứ nhất là $2$, thừa số thứ hai là $9$. Tích của hai số đó là
\choice
{$16$}
{$17$}
{$20$}
{\True $18$}
\loigiai{
$2 \times 9 = 18$.
\textbf{Đáp án của câu hỏi là D}
}
\end{ex}

\begin{ex} %[2D8V38]
Trong phép nhân có thừa số là $5$, nếu tăng thừa số thứ hai thêm $2$ đơn vị thì tích tăng thêm bao nhiêu?
\choice
{\True $10$ đơn vị}
{$5$ đơn vị}
{$7$ đơn vị}
{$2$ đơn vị}
\loigiai{
Thừa số thứ hai tăng 2 đơn vị thì tích tăng thêm: $5 \times 2 = 10$ đơn vị.
\textbf{Đáp án của câu hỏi là A}
}
\end{ex}

\begin{ex} %[2D8V38]
Tìm một thừa số biết tích bằng $14$ và thừa số kia bằng $2$.
\choice
{$6$}
{$8$}
{\True $7$}
{$9$}
\loigiai{
$14 : 2 = 7$.
\textbf{Đáp án của câu hỏi là C}
}
\end{ex}

\begin{ex} %[2D8V38]
Tích của số chẵn nhỏ nhất có một chữ số khác 0 và số lớn nhất có một chữ số là
\choice
{$16$}
{\True $18$}
{$14$}
{$20$}
\loigiai{
Số chẵn nhỏ nhất khác 0 là 2. Số lớn nhất có 1 chữ số là 9. Tích: $2 \times 9 = 18$.
\textbf{Đáp án của câu hỏi là B}
}
\end{ex}

\begin{ex} %[2D8V38]
Điền số thích hợp: $2 \times \dots = 16$.
\choice
{$6$}
{$7$}
{$9$}
{\True $8$}
\loigiai{
$16 : 2 = 8$.
\textbf{Đáp án của câu hỏi là D}
}
\end{ex}

\begin{ex} %[2D8V38]
Có hai túi kẹo, mỗi túi có $6$ cái kẹo. Hỏi cả hai túi có bao nhiêu cái kẹo?
\choice
{\True $12$ cái kẹo}
{$10$ cái kẹo}
{$14$ cái kẹo}
{$8$ cái kẹo}
\loigiai{
$6 \times 2 = 12$ (cái kẹo).
\textbf{Đáp án của câu hỏi là A}
}
\end{ex}

\begin{ex} %[2D8V38]
Tích của $5$ với số nào cho kết quả bằng $35$?
\choice
{$6$}
{$8$}
{\True $7$}
{$9$}
\loigiai{
$35 : 5 = 7$.
\textbf{Đáp án của câu hỏi là C}
}
\end{ex}

\begin{ex} %[2D8V38]
Trong một phép nhân, nếu một thừa số bằng $0$ thì tích luôn bằng
\choice
{$1$}
{\True $0$}
{Chính thừa số kia}
{Không xác định}
\loigiai{
Bất kỳ số nào nhân với 0 cũng bằng 0.
\textbf{Đáp án của câu hỏi là B}
}
\end{ex}

\begin{ex} %[2D8V38]
Tính giá trị biểu thức: $2 \times 8 - 7$.
\choice
{$8$}
{$10$}
{$7$}
{\True $9$}
\loigiai{
$2 \times 8 = 16$; $16 - 7 = 9$.
\textbf{Đáp án của câu hỏi là D}
}
\end{ex}

\begin{ex} %[2D8V38]
Tìm tích của số liền sau số $1$ và số liền trước số $6$.
\choice
{\True $10$}
{$8$}
{$12$}
{$9$}
\loigiai{
Số liền sau 1 là 2. Số liền trước 6 là 5. Tích: $2 \times 5 = 10$.
\textbf{Đáp án của câu hỏi là A}
}
\end{ex}

\begin{ex} %[2D8V38]
Mỗi bạn nhỏ có $2$ quả bóng bay. Hỏi $7$ bạn nhỏ có bao nhiêu quả bóng bay?
\choice
{$12$ quả}
{$16$ quả}
{\True $14$ quả}
{$9$ quả}
\loigiai{
$2 \times 7 = 14$ (quả).
\textbf{Đáp án của câu hỏi là C}
}
\end{ex}

\textbf{Phần 2. Trắc nghiệm đúng sai}
\setcounterpageref{ex}{0}

\begin{ex} %[2D8V38]
Xét tính đúng sai của các nhận định về thừa số và tích:
\choiceTF
{\True Trong phép nhân a x b = c, a và b là các thừa số, c là tích.}
{Đổi chỗ hai thừa số thì tích sẽ bị thay đổi.}
{\True Tích của 2 số tự nhiên khác 0 luôn lớn hơn hoặc bằng mỗi thừa số.}
{Nếu tích bằng 0 thì cả hai thừa số đều bắt buộc phải bằng 0.}
\loigiai{
\begin{itemize}
    \item Ý a) Đúng: Đúng theo tên gọi.
    \item Ý b) Sai: Tính chất giao hoán: đổi chỗ hai thừa số thì tích không đổi.
    \item Ý c) Đúng: Đúng.
    \item Ý d) Sai: Chỉ cần ít nhất một thừa số bằng 0 là tích bằng 0.
\end{itemize}
\begin{itemize}
    \item a) \textbf{Đúng}.
    \item b) \textbf{Sai}.
    \item c) \textbf{Đúng}.
    \item d) \textbf{Sai}.
\end{itemize}
}
\end{ex}

\begin{ex} %[2D8V38]
Cho phép nhân 5 x 6 = 30:
\choiceTF
{\True Các thừa số là 5 và 6, tích là 30.}
{\True Nếu tăng thừa số thứ nhất thêm 1 đơn vị thành 6 thì tích mới là 36.}
{Tích 30 là một số lẻ.}
{\True Biểu thức 5 x 6 bằng biểu thức 6 x 5.}
\loigiai{
\begin{itemize}
    \item Ý a) Đúng: Đúng.
    \item Ý b) Đúng: $6 \times 6 = 36$. Đúng.
    \item Ý c) Sai: 30 có chữ số tận cùng là 0 nên là số chẵn.
    \item Ý d) Đúng: Đúng.
\end{itemize}
\begin{itemize}
    \item a) \textbf{Đúng}.
    \item b) \textbf{Đúng}.
    \item c) \textbf{Sai}.
    \item d) \textbf{Đúng}.
\end{itemize}
}
\end{ex}

\textbf{Phần 3. Trả lời ngắn}
\setcounterpageref{ex}{0}

\begin{ex} %[2D8V38]
Tìm tích của 2 và 7.
\shortans{14}
\loigiai{
\textbf{Đáp số:} $14$\\
$2 \times 7 = 14$.
}
\end{ex}

\begin{ex} %[2D8V38]
Tìm thừa số chưa biết: x x 5 = 40.
\shortans{8}
\loigiai{
\textbf{Đáp số:} $8$\\
$x = 40 : 5 = 8$.
}
\end{ex}

\begin{ex} %[2D8V38]
Tính giá trị: 2 x 9 + 15.
\shortans{33}
\loigiai{
\textbf{Đáp số:} $33$\\
$2 \times 9 = 18$; $18 + 15 = 33$.
}
\end{ex}

\begin{ex} %[2D8V38]
Mỗi nhóm học sinh có 5 bạn. Hỏi 6 nhóm như thế có tất cả bao nhiêu bạn học sinh?
\shortans{30}
\loigiai{
\textbf{Đáp số:} $30$\\
$5 \times 6 = 30$ bạn.
}
\end{ex}

\begin{ex} %[2D8V38]
Tìm tích của số lớn nhất có một chữ số và số 2.
\shortans{18}
\loigiai{
\textbf{Đáp số:} $18$\\
$9 \times 2 = 18$.
}
\end{ex}

\begin{ex} %[2D8V38]
Một lọ hoa cắm 5 bông hoa. Hỏi 8 lọ hoa như thế có bao nhiêu bông hoa?
\shortans{40}
\loigiai{
\textbf{Đáp số:} $40$\\
$5 \times 8 = 40$ bông hoa.
}
\end{ex}
